\mxunit[5]{terminale-c05}{Combinatoire et dénombrement}\label{ch:9}

\mxsubtitle{S22, S25, S30, S35}

\begin{keypoint}[Objectifs]Choisir un modèle de comptage ; distinguer ordre et répétition ; utiliser coefficients binomiaux et relation de Pascal.\end{keypoint}

\section{Notes de cours}
\subsection{Additionner ou multiplier}
Pour des ensembles finis deux à deux disjoints $A_i$, $|\bigcup A_i|=\sum|A_i|$. Si les cas se recouvrent, il faut corriger le double comptage : $|A\cup B|=|A|+|B|-|A\cap B|$.
Le nombre d'éléments du produit cartésien $A\times B$ est $|A||B|$. Pour une suite de choix, multiplier les nombres de possibilités à chaque étape, à condition que ces nombres soient valables pour chaque choix antérieur considéré.
\paragraph{Exemple.} Un code de deux lettres choisies parmi $26$, suivies de trois chiffres, avec répétition autorisée, peut prendre $26^2\times10^3$ valeurs. Un nombre à trois chiffres ne commence pas par $0$ : il y en a $9\times10^2$.
\subsection{Listes et parties}
Dans un ensemble de $n$ éléments, il y a $n^k$ listes ordonnées de longueur $k$ avec répétition ; $n(n-1)\cdots(n-k+1)=\frac{n!}{(n-k)!}$ listes de $k$ éléments distincts, si $0\le k\le n$ ; $n!$ permutations de tous les éléments. On pose $0!=1$.
\begin{keypoint}[Combinaisons]
Une combinaison de $k$ éléments parmi $n$ est une partie à $k$ éléments ; l'ordre n'intervient pas.
\[\binom nk=\frac{n!}{k!(n-k)!},\quad\binom n0=\binom nn=1,\quad\binom n1=n,\quad\binom n2=\frac{n(n-1)}2.\]
\end{keypoint}
\paragraph{Pourquoi diviser par $k!$ ?} Chaque partie de taille $k$ donne exactement $k!$ listes distinctes selon l'ordre de ses éléments. Le nombre de listes sans répétition est donc $k!\binom nk$.
\subsection{Identités et démonstrations}
\[\binom nk=\binom n{n-k},\qquad \binom nk+\binom n{k+1}=\binom{n+1}{k+1},\qquad\sum_{k=0}^n\binom nk=2^n.\]
La relation de Pascal ci-dessus s'utilise pour $0\le k<n$ ; les formules au bord se lisent directement avec les valeurs $1$.
\paragraph{Symétrie.} Prendre le complémentaire établit une correspondance entre les parties de taille $k$ et celles de taille $n-k$.
\paragraph{Pascal, preuve combinatoire.} Dans un ensemble de $n+1$ éléments, fixer un élément $a$. Une partie de taille $k+1$ contient $a$ (choisir $k$ autres éléments parmi $n$) ou ne le contient pas (choisir $k+1$ éléments parmi $n$). Les deux cas sont disjoints et couvrent toutes les possibilités.
\paragraph{Pascal, preuve algébrique.} Pour $0\le k<n$, mettre au même dénominateur :
\[\frac{n!}{k!(n-k)!}+\frac{n!}{(k+1)!(n-k-1)!}
=\frac{n!((k+1)+(n-k))}{(k+1)!(n-k)!}
=\frac{(n+1)!}{(k+1)!(n-k)!}.\]
Les cas au bord se traitent directement.
\paragraph{Nombre de parties.} Pour chacun des $n$ éléments, choisir « présent » ou « absent » donne $2^n$ parties. En les classant par taille, on en compte $\sum_{k=0}^n\binom nk$ : les deux comptages sont égaux.
\paragraph{Approfondissement mentionné dans la progression.} Le développement $(a+b)^n=\sum_{k=0}^n\binom nk a^{n-k}b^k$ se comprend en choisissant, dans chacun des $n$ facteurs, le terme $a$ ou $b$. Cette formule n'est pas utilisée comme prérequis des DS de ce livret.


\mxnewpage
\section{Activité de recherche}\textbf{Ordre et répétition dans un dénombrement --- 30 min.}\par Six élèves distincts se portent volontaires pour organiser un débat.
\begin{enumerate}\item On choisit deux élèves pour former un binôme sans attribuer de rôle. Combien de binômes peut-on former ?
\item On désigne un président et un secrétaire, qui doivent être deux élèves différents. Combien d'attributions des deux rôles sont possibles ?
\item On prévoit trois interventions successives. Un même élève peut intervenir plusieurs fois. Combien de programmes de trois interventions sont possibles ?
\item On recommence avec quatre élèves seulement, désignés par $A$, $B$, $C$ et $D$. Écrire tous les binômes et toutes les attributions des deux rôles. Expliquer ensuite pourquoi il existe $64$ programmes de trois interventions, sans les énumérer.
\item Pour chacun des trois modes de choix, préciser si l'ordre intervient et si une même personne peut être choisie plusieurs fois.\end{enumerate}

\mxnewpage
\section{Exercices progressifs}

\mxexerciseitem{1}{Entraînement}Un code est une suite de quatre chiffres choisis parmi $0,1,\ldots,9$. Les répétitions et le chiffre $0$ en première position sont autorisés. Combien de codes contiennent au moins un $0$ ?

\mxexerciseitem{2}{Entraînement}Dix athlètes participent à une course sans ex æquo. Combien de podiums différents, avec une première, une deuxième et une troisième place, peut-on obtenir ?

\mxexerciseitem{3}{Entraînement}Parmi sept personnes distinctes, on forme un groupe de trois personnes contenant une personne désignée à l'avance. Combien de groupes peut-on former ?

\mxexerciseitem{4}{Entraînement}Sur un quadrillage, on part de $\pt{0;0}$ pour atteindre $\pt{4;3}$ en effectuant exactement quatre déplacements d'une unité vers la droite et trois vers le haut. Combien de chemins différents peut-on suivre ?

\mxexerciseitem{5}{Entraînement}Une classe compte huit filles et sept garçons. Combien de groupes de quatre élèves contenant exactement deux filles peut-on former ?

\section{Travail dans le manuel ou une annale}\textbf{Barbazo, p. 294, exercice 37.} À la maison après S35 : traiter les deux questions. Préciser si l'ordre intervient et si une carte peut être choisie plusieurs fois avant d'utiliser un coefficient binomial.

\mxnewpage
\section{Devoir à la maison}\textbf{DM9 --- Former une commission --- 45 min, après S35.}\par Une équipe comprend cinq enseignants et sept élèves, tous distincts. On forme une commission de quatre personnes.
\begin{enumerate}\item Combien de commissions différentes peut-on former ?
\item Combien de commissions contiennent exactement deux élèves ?
\item Combien de commissions contiennent au moins un enseignant et au moins un élève ?
\item Parmi les membres d'une commission de la question 2, on désigne un président enseignant et un secrétaire élève. Dénombrer les choix de la commission accompagnée de ces deux attributions.
\item En dénombrant les commissions selon leur nombre d'enseignants, démontrer que \[\sum_{k=0}^4\binom5k\binom7{4-k}=\binom{12}4.\]\end{enumerate}

\ifmxteacher\mxnewpage\section{Conduite et corrigés}\begin{teachernote}[Repères pour le professeur]Le devoir commun S28 précède les combinaisons de S30 : n'y mettre que principes et listes ordonnées. Le DS3 peut mobiliser les combinaisons, mais pas la relation de Pascal enseignée en S35. Faire formuler l'objet compté avant toute formule. Pour les élèves fragiles, utiliser des cartes portant des noms ; une permutation de cartes change-t-elle l'objet ? Le binôme de Newton, cité dans la progression, est un approfondissement : le garder explicitement facultatif.\end{teachernote}

\subsection{Corrigé de l'activité}\begin{enumerate}
\item Un binôme est une partie de deux élèves parmi six : $\binom{6}{2}=\frac{6\times5}{2}=15$.
\item Le président peut être choisi de six façons, puis le secrétaire parmi les cinq autres élèves : $6\times5=30$ répartitions.
\item Chaque intervention peut être confiée à l'un des six élèves, y compris à un élève déjà choisi. Les trois choix successifs donnent $6^3=216$ suites.
\item Avec quatre cartes, les binômes sont $AB,AC,AD,BC,BD,CD$ : six choix. Les douze couples de rôles sont :\[\begin{array}{lll}(A,B)&(A,C)&(A,D)\\(B,A)&(B,C)&(B,D)\\(C,A)&(C,B)&(C,D)\\(D,A)&(D,B)&(D,C)\end{array}\] Les suites de trois interventions sont au nombre de $4^3=64$, car chaque position possède quatre choix.
\item Pour un binôme, l'ordre ne compte pas et on ne répète pas un élève. Pour les rôles, l'ordre compte, sans répétition. Pour les interventions, l'ordre compte et la répétition est autorisée.
\end{enumerate}

\paragraph{E1.}Il existe $10^4$ codes au total. Ceux qui ne comportent aucun $0$ ont neuf choix à chacune des quatre positions, soit $9^4$ codes. Par complément, le nombre demandé est $10^4-9^4=3439$.

\paragraph{E2.}Il y a dix choix pour l'or, puis neuf pour l'argent et huit pour le bronze. Le principe multiplicatif donne $10\times9\times8=720$ podiums. Changer l'ordre change le podium.

\paragraph{E3.}La personne imposée est déjà choisie. Il reste à choisir une partie de deux personnes parmi les six autres : $\binom{6}{2}=\frac{6\times5}{2}=15$.

\paragraph{E4.}Un chemin est déterminé par les positions de ses quatre pas à droite parmi sept positions. Les trois autres pas sont alors imposés. On compte donc $\binom74=35$ chemins.

\paragraph{E5.}Il faut choisir deux filles parmi huit, puis deux garçons parmi sept. Chaque couple de choix donne un seul groupe : $\binom82\binom72=28\times21=588$.

\mxneed{30}

\subsection{Corrigé du DM9}\begin{enumerate}
\item L'ordre ne compte pas : il y a $\binom{12}{4}=\frac{12\times11\times10\times9}{24}=495$ commissions.
\item On choisit deux élèves parmi sept et deux enseignants parmi cinq : $\binom72\binom52=21\times10=210$ commissions.
\item On retire des $495$ commissions les $\binom54=5$ commissions exclusivement enseignantes et les $\binom74=35$ commissions exclusivement élèves. Ces catégories sont disjointes : $495-5-35=455$.
\item Chacune des $210$ commissions possède deux choix de président enseignant et deux choix de secrétaire élève : $210\times2\times2=840$ organisations.
\item Classons les commissions selon le nombre $k$ d'enseignants, avec $0\le k\le4$. La catégorie $k$ contient $\binom5k\binom7{4-k}$ commissions. Les cinq catégories sont disjointes et couvrent toutes les commissions. Par addition,
\[\sum_{k=0}^4\binom5k\binom7{4-k}=\binom{12}4=495.\]
\end{enumerate}

\subsection{Corrigé du travail dans le manuel ou une annale}Une main contient cinq cartes distinctes ; l'ordre n'intervient pas. Il y a donc $\binom{52}5=\num{2598960}$ mains. Pour une main contenant l'as de pique, cette carte est déjà choisie ; on choisit les quatre autres parmi les $51$ cartes restantes. Il y a $\binom{51}4=\num{249900}$ mains. On ne multiplie pas ce résultat par $5$ : la main n'a pas de première ou de dernière position.

\fi

